%!TEX root = ../MainBody.tex

% 中英文摘要
%% ═══════════════════════════════════════════════════════════════════
%% [待办] 以下为模板占位文字，需全部替换为论文实际摘要。
%%
%% 中文摘要写作思路（约 800-1000 字）：
%%   第 1 段：研究背景与问题
%%     - 红外仿真的重要性、应用领域与物理复杂性
%%     - 简要介绍FEM方法。
%%     - 基于路径追踪的蒙特卡洛耦合热输运方法的优势与计算瓶颈
%%     - 现有 CPU DFS 架构的局限性
%%   第 2 段：研究内容与方法
%%     - 贡献 1：Wavefront 耦合热输运求解器（DFS→状态机→Wavefront 并行化）
%%     - 贡献 2：基于 OptiX 的 GPU 硬件光追查询后端（三类批量查询）
%%     - 贡献 3：GUI 场景编辑器（场景编辑/可视化/任务管理）
%%   第 3 段：实验结果
%%     - 物理一致性验证（GPU/CPU 容差 1e-6）
%%     - 性能加速比（需填入最终数据）
%%   第 4 段：结论性陈述
%%     - 本文方法的意义和应用价值
%%
%% [数据需求] 最终加速比数字（待实验完成后填写）
%%
%% 英文摘要：中文摘要的对应翻译，结构一致
%% ═══════════════════════════════════════════════════════════════════
\begin{CHSabstract}
红外辐射携带目标温度分布与表面辐射特性信息，在目标探测识别、精确制导、工业无损检测与医疗热成像等领域具有重要应用价值。红外仿真通过计算机预测场景的红外辐射特性，是装备研制、性能评估与对抗仿真的重要支撑手段。物理自洽的红外仿真须建模并求解由辐射、传导、对流及其在界面上耦合构成的耦合热输运过程，该过程由一组非线性偏微分方程描述。

该方程组的求解一般有有限元方法与蒙特卡洛方法两种途径。有限元方法对方程组进行空间离散，以有限维线性方程组近似各物理场分布并迭代求解，多过程耦合在迭代外联立。但该方法精度受迭代策略和网格质量制约，自由度总数随场景规模急剧膨胀。蒙特卡洛方法对方程组积分形式进行统计估计，收敛速度无关于问题维度，不依赖空间离散，适应复杂边界条件与多介质耦合界面。近年来，基于路径追踪的耦合求解框架将辐射、传导与对流纳入单一路径空间；然而，已有实现采用递归路径追踪模式，在结构上阻断了跨路径积累同类查询以形成批次的可能，从而难以利用 GPU 专用光线追踪单元的并行能力，成为制约该方法进一步规模化的瓶颈。

本文针对上述瓶颈提出 Wavefront 架构耦合热输运求解框架，包含三项研究内容。

第一，提出基于有限状态机的 Wavefront 求解算法。将递归算法中的几何查询定义为挂起点，围绕挂起点推导出路径状态机的完备设计约束、数据结构与状态转换规则；经递归路径追踪算法到有限状态机的等价变换，将路径中间状态提取为持久化路径状态结构体，解耦状态推进与查询执行；以持久化路径池为中心构建 Wavefront 迭代执行模型，跨路径积累查询后批量提交，实现大规模并行几何查询。

第二，提出三类异构几何查询的 GPU 等价映射方案。射线最近命中查询由专用光追单元直接加速；场景最近点查询通过几何等价重构转化为等价光追查询形式，使其同样满足硬件加速的算法结构；封闭腔体成员判定查询基于 Jordan--Brouwer 分离定理，分解为上述两类查询的组合，以两次批量查询完成判定；三类映射整合于统一 GPU 查询管线，从而使耦合热输运求解过程中的全部几何查询均可利用 GPU 专用光线追踪单元加速。

第三，作为面向仿真工作流的配套工程工具，以求解器输入语言 \texttt{stardis-input} 的 BNF 形式文法为基础，设计并实现图形化场景编辑器 StardisEditor。该 DSL 采用严格位置编码的纯文本格式，存在位置参数无标签、几何关系不可见、表面区域划分依赖外部建模工具等系统性使用障碍；编辑器以 BNF 文法规范驱动语法感知的解析与写出闭环，在此基础上提供三项 DSL 能力扩展——基于 VTK 的三维几何可视化、交互式画笔标注表面区域、图形化任务配置与队列化执行——将仿真建模、验证与任务管理整合于统一图形化工作流。

实验验证表明：所提出的Wavefront求解算法在涵盖传导、辐射、对流及多机制耦合的13个标准测试场景的全部282个数据点上，均通过$3\sigma$检验，逐像素统计检验经多重比较校正后报告零显著差异，证明算法等价变换不引入系统性物理偏差。在 8\textasciitilde 16 线程典型工作站配置下，Wavefront 求解器配合 GPU 光追后端实现 $1.3$\textasciitilde$1.4$ 倍端到端加速，证明了 Wavefront 架构的价值和可行性。图形化编辑器通过 Cornell Box 耦合热仿真案例验证了从场景构建到红外渲染输出的完整闭环工作流。
\end{CHSabstract}

\begin{ENGabstract}
Infrared radiation encodes the temperature distribution and emissive properties of objects and has important applications in target detection and recognition, precision guidance, industrial non-destructive inspection, and medical thermal imaging. Infrared simulation---computationally predicting the infrared signature of a scene---is an essential tool for system design, performance evaluation, and countermeasure assessment. Physically self-consistent infrared simulation requires modeling and solving the coupled heat transport process comprising radiation, conduction, convection, and their interfacial coupling, which is governed by a system of nonlinear partial differential equations.

Two general approaches exist for solving these equations. The finite element method spatially discretizes the equations, approximating each physical field by a finite-dimensional linear system solved iteratively, with multi-process coupling handled through outer iteration loops. However, its accuracy is constrained by the iterative strategy and mesh quality, and the total degrees of freedom grow rapidly with scene complexity. The Monte Carlo method statistically estimates the integral form of the equations, with a convergence rate independent of problem dimensionality, requiring no spatial discretization and adapting naturally to complex boundary conditions and multi-material coupling interfaces. Recently, a path-tracing-based coupled solution framework has unified radiation, conduction, and convection within a single path space; however, existing implementations employ a recursive path tracing model that structurally prevents homogeneous queries from being accumulated across paths into batches, precluding effective use of the parallel capability of dedicated GPU ray tracing units and forming the central bottleneck limiting further scalability.

This thesis addresses the above bottleneck by proposing a wavefront coupled heat transport solution framework, comprising three contributions.

First, a wavefront solving algorithm based on finite state machines is proposed. Geometry query call sites in the recursive algorithm are defined as suspension points; complete design constraints, data structures, and state transition rules for the path state machine are derived around these suspension points. Through an equivalent transformation from recursive path tracing to a finite state machine, path intermediate states are lifted into persistent path state structures, decoupling state advancement from query execution. A wavefront iterative execution model centered on a persistent path pool accumulates queries across paths and submits them in batches, enabling large-scale parallel geometry processing.

Second, GPU-equivalent mappings for three heterogeneous geometry query types are derived. Ray nearest-hit queries are directly accelerated by dedicated GPU ray tracing units. Closest-point queries are reformulated into an equivalent ray tracing form through geometric equivalence, satisfying the algorithmic structure required for hardware acceleration. Enclosed cavity membership queries are decomposed via the Jordan--Brouwer separation theorem into a combination of the above two query types, completing the determination in two batch invocations. All three mappings are integrated into a unified GPU query pipeline, enabling every geometry query in the coupled heat transport solving process to be accelerated by dedicated GPU ray tracing units.

Third, as an engineering companion tool for the simulation workflow, a graphical scene editor, StardisEditor, is designed and implemented on the basis of the BNF formal grammar of the solver's input language \texttt{stardis-input}. This DSL adopts a strictly position-encoded plain-text format, giving rise to systematic usability barriers including unlabeled positional parameters, invisible geometric relationships, and reliance on external modeling tools for surface-region partitioning. The editor employs a grammar-aware parsing and serialization loop driven by the BNF specification and extends the DSL's capabilities in three directions: VTK-based three-dimensional geometry visualization, interactive brush-based surface-region labeling, and graphical task configuration with queued execution, unifying scene modeling, validation, and task management within a single graphical workflow.

Experimental validation shows that the wavefront solving algorithm passes the $3\sigma$ criterion on all 282 data points across 13 test scenes, and per-pixel statistical tests report zero significant differences after multiple comparison correction, confirming no systematic physical bias is introduced by the equivalent transformation. Under typical 8--16 thread workstation configurations, the wavefront solver combined with the GPU ray tracing backend achieves $1.3$--$1.4\times$ end-to-end speedup, demonstrating the value and feasibility of the wavefront architecture. The graphical editor is validated through a Cornell Box coupled thermal simulation case that exercises the complete closed-loop workflow from scene construction to infrared rendering output.
\end{ENGabstract}

